\documentclass[11pt]{article}

\usepackage[margin=1in]{geometry}
\usepackage{booktabs}
\usepackage{amsmath}

\setlength{\emergencystretch}{2em}

\title{LLM-Assisted Entity Resolution:\\Summary of Findings and Current Status}
\author{}
\date{2026-08-02}

\begin{document}
\maketitle

\section{Purpose}

This report describes how the entity-resolution pipeline works, explains each
attribute-selection methodology it implements, and summarizes what the experiments so
far establish. The methodology sections are complete; the results are condensed. Full
per-run provenance and the detailed statistical analysis behind each finding are in
the companion report dated 2026-08-01.

All numbers come from runs matching the pipeline's current configuration: embedding
model \texttt{BAAI/bge-large-en-v1.5}, matcher LLM \texttt{google/gemini-2.5-flash} via
OpenRouter. Fodors--Zagat runs exclude the leaking \texttt{class} attribute
(Section~\ref{sec:leak}); runs whose LLM calls failed are excluded entirely
(Section~\ref{sec:defects}).

\section{Pipeline Overview}

The pipeline runs in five stages, each timed and recorded in a per-run JSON summary
under \texttt{logs/<dataset>/runs/}.

\textbf{Load and normalize.} Both tables and the ground truth are read, ID columns
are dropped, and every remaining field is passed through a normalization step that
expands or standardizes address and directional abbreviations, and cleans up
Unicode/HTML artifacts, so that superficial formatting differences do not by
themselves create false mismatches later in the pipeline.

\textbf{Embed and block.} Each attribute column is encoded independently with the
SentenceBERT model, the column vectors are concatenated and L2-normalized into one
embedding per record, and multi-probe locality-sensitive hashing (random
hyperplanes, with single-bit flips at query time) retrieves the top-$k$ most similar
Table-B candidates for every Table-A record. This reduces the pairwise comparison
space from the full cross product down to a much smaller candidate-pair set before
any LLM is involved.

\textbf{Step 1 -- global tuple condensation (optional).} One attribute subset is
chosen once for the whole dataset and applied to every candidate pair before
re-blocking on the condensed attributes. This step is described in full in
Section~\ref{sec:step1}.

\textbf{Step 1.5 -- adaptive per-pair attribute selection (optional).} Instead of one
global subset, a small sample of candidate pairs is profiled by the LLM for
per-attribute importance, and a hybrid predictor/retriever model transfers a
\emph{different} attribute subset to every candidate pair. This step is described in
full in Section~\ref{sec:step15}.

\textbf{Match.} One LLM call decides Yes/No per candidate pair, using only the
attributes selected for that pair (all attributes if neither Step 1 nor Step 1.5
ran). Every response is cached on disk keyed by a hash of the model, prompt
template, and record content, so repeated experiments over the same configuration
do not re-spend tokens.

\textbf{Metrics and logging.} Precision, recall, and F1 are computed both at the
candidate-pair level (only pairs that survived blocking) and end-to-end (relative to
all ground-truth matches, so blocking misses count against recall), alongside
token-usage accounting split by stage.

\section{Methodology}
\label{sec:methodology}

\subsection{Step 1: Global Tuple Condensation}
\label{sec:step1}

The motivation for Step 1 is that showing an LLM every attribute for every pair
wastes tokens on fields that are frequently empty, redundant, or simply irrelevant
to the match decision, and can make the model's decision worse by diluting the
signal from the few attributes that actually distinguish entities. Step 1 picks
\emph{one} attribute subset for the whole dataset, used for every pair, via one of
four strategies.

\paragraph{Manual selection.} A fixed, user-supplied list of attribute names is kept
and everything else is dropped, with no data-driven reasoning involved. It is the
cheapest strategy to run (no LLM calls, no labeled data required) and is appropriate
when a person can already tell from the schema which fields carry identity
information, but it does not adapt if that judgment is wrong for some records.

\paragraph{Heuristic selection.} Attributes are pruned using two label-free,
corpus-level statistics: the fraction of missing values (an attribute that is empty
in most records cannot be evidence of anything) and the fraction of unique values
(an attribute that takes almost the same value everywhere carries little
discriminative signal). Attributes that survive pruning are also \emph{compressed}:
long text fields are stripped of boilerplate words (via a stopword/marketing-language
regex list) and only high-signal tokens are kept -- numbers, alphanumeric
codes/identifiers, capitalized proper nouns, and unit-bearing measurements (e.g.\
\texttt{"12 kg"}) -- so that a lengthy description field still contributes its
distinctive tokens without its full token count.

\paragraph{LLM-guided selection.} A small labeled sample of matching and
non-matching pairs is shown to the LLM one pair at a time, together with the
question ``which fields were most influential in deciding whether these two records
match?'', and the model returns the smallest sufficient subset of attribute names for
that pair. Responses are aggregated across the sample by \emph{selection frequency}
per attribute (what fraction of responses named that attribute), attributes are
ranked by that frequency, and the final subset keeps every attribute above a
frequency threshold (optionally capped to a top-$k$).

\paragraph{Supervised selection.} For the same labeled sample, three field-level
similarity features are computed per attribute for every pair -- token-set Jaccard
overlap, normalized edit-distance similarity, and an exact-match indicator -- and a
logistic regression classifier (features standardized, class-balanced) is trained to
predict match/non-match from all attributes' features at once. Each attribute's
importance is the maximum absolute regression weight across its three features,
normalized so all attributes' importance scores sum to one; the top-ranked
attributes above a threshold (optionally capped to top-$k$) are kept. Because it
learns from labeled evidence directly, this strategy tends to identify a very small,
highly discriminative attribute subset (see Section~\ref{sec:results}).

Whichever strategy is used, the condensed table is then re-embedded and re-blocked
(so blocking itself also benefits from the smaller, more focused attribute set)
before matching.

\subsection{Step 1.5: Pair-Adaptive Attribute Selection}
\label{sec:step15}

Step 1.5 addresses a limitation of Step 1: even a well-chosen global attribute
subset is not equally informative for every pair. For one pair an exact model
number may already be decisive; for another, the deciding signal might be an author
mismatch or a subtle title variation, while the model number field is missing or
uninformative. Step 1.5 replaces one global subset with a \emph{per-pair} subset,
produced in two stages.

\textbf{Profiling stage.} A small sample of post-blocking candidate pairs (the
``profile pairs'') is selected \emph{without using ground-truth labels} by one of
five samplers (described below). For each profile pair, the LLM is asked to score
every attribute's usefulness for the match/non-match decision on an integer 0--3
scale (0 = not useful/missing/generic, 3 = highly important or decisive), and the
raw scores are normalized per pair into an importance vector that sums to one.

\textbf{Transfer stage.} A \texttt{HybridAttributeSelector} is fit on the profiled
pairs' label-free similarity features (token overlap, edit similarity, exact match,
semantic similarity, numeric similarity, and an is-numeric indicator, computed per
attribute) against their LLM-scored importance vectors. It combines two
sub-models: a Ridge-regression \emph{predictor} that maps a pair's features directly
to a predicted importance vector, and a cosine-similarity \emph{retriever} that finds
the $k$ most similar profiled pairs and averages their importance vectors. The two
signals are combined by a weighted average (0.5/0.5 in the current configuration).
For every candidate pair (profiled or not), attributes are then selected in
descending combined-importance order until either a fixed top-$k$ is reached or
cumulative importance crosses a threshold (0.8 in the current configuration,
bounded between a minimum and maximum attribute count), with a fallback rule that
force-includes at least one attribute tagged with a required ``identity'' role
(e.g.\ title, name, model number) even if the learned ranking would otherwise drop
it.

The five samplers below differ only in \emph{which} candidate pairs get profiled by
the LLM; the transfer stage is identical for all of them.

\paragraph{Random sampling.} Profile pairs are drawn uniformly at random from all
candidate pairs. This is the label-free baseline against which the other samplers
are compared.

\paragraph{Similarity-stratified sampling.} Candidate pairs are bucketed into
low/medium/high terciles by their blocking-stage cosine similarity, and profile
pairs are drawn from each bucket in fixed proportions (60\% medium, 20\% low, 20\%
high in the current configuration). This intentionally profiles some clearly
similar pairs, some clearly dissimilar pairs, and a majority of ambiguous
near-threshold pairs, on the reasoning that ambiguous pairs are where attribute
importance is hardest to guess and most useful to learn from.

\paragraph{Diversity sampling.} Profile pairs are chosen by farthest-first traversal
in the label-free pair-feature space: starting from the pair farthest from the
sample centroid, each subsequent pair is the one that maximizes its minimum
distance to all previously chosen pairs. This favors broad coverage of distinct
similarity ``shapes'' (e.g.\ pairs that agree on names but disagree on numbers, versus
pairs that disagree on everything) over any particular similarity level.

\paragraph{Stratified diversity sampling.} The same low/medium/high similarity
strata as similarity-stratified sampling are used, but within each stratum, pairs
are chosen by the same farthest-first diversity rule rather than uniformly at
random. This is meant to combine both benefits: similarity-level balance and
within-level feature diversity.

\paragraph{Importance-based sampling.} Pairs are ranked by a label-free
\emph{importance proxy score} that estimates, without ground truth, how
informative a pair is likely to be for learning attribute importance:
$$\text{proxy} = 0.50 \cdot \text{evidence\_contrast} + 0.30 \cdot \text{similarity\_uncertainty} + 0.20 \cdot \text{evidence\_mixedness}$$
%
where evidence\_contrast is the standard deviation of per-attribute agreement
signals (favoring pairs where attributes disagree with each other, rather than all
pointing the same way), similarity\_uncertainty favors pairs whose blocking
similarity is far from the sample's typical value, and evidence\_mixedness favors
pairs whose average per-attribute agreement sits near 0.5 (neither clearly matching
nor clearly non-matching). The top-scoring pairs (with random tie-breaking) become
the profile sample.

\section{Results and Findings}
\label{sec:results}

\subsection{Attribute selection is worth doing}

On the one dataset with a no-selection baseline, both mechanisms beat showing the
matcher every field, and both spend fewer tokens doing it. This is the premise
underlying the whole approach, and it holds.

\begin{center}
\small
\begin{tabular}{lrrrr}
\toprule
Method & Attrs/pair & F1 & Tokens & vs.\ no selection \\
\midrule
No selection (\texttt{full}) & 5.00 & 0.7849 & 802{,}477 & --- \\
Step 1.5, diversity (mean of 6 runs) & 4.00 & 0.9222 & 775{,}828 & $+0.1373$ F1, $-3.3\%$ tokens \\
Step 1, supervised & 2.00 & \textbf{0.9537} & \textbf{658{,}185} & $+0.1688$ F1, $-18.0\%$ tokens \\
\bottomrule
\end{tabular}
\end{center}

\noindent Fodors--Zagat, leaking attribute excluded. Note that the simpler mechanism
wins here: a single global two-attribute subset beats per-pair adaptive selection on
both accuracy and cost. Section~\ref{sec:supervision} qualifies what that comparison
actually establishes.

\subsection{The two mechanisms do not operate under the same supervision}
\label{sec:supervision}

The head-to-head above is not like-for-like, and the difference is large enough to
change how the result should be read.

Step 1's supervised strategy requires ground-truth labels. It calls
\texttt{build\_step1\_labeled\_pairs}, which draws its positive examples directly from
the gold set --- 50 positives and 50 negatives in the configuration used here --- and
trains a classifier on them. Step 1's LLM-guided strategy needs the same labeled
sample. Step 1.5, by contrast, consumes \emph{no} labels at all: its profile pairs are
chosen by label-free samplers and scored by the LLM, never by consulting ground truth.

\begin{center}
\small
\begin{tabular}{lrrr}
\toprule
Mechanism & Gold labels used & Fodors--Zagat & DBLP--ACM \\
\midrule
Step 1, supervised / LLM-guided & 50 pos.\ $+$ 50 neg. & 45.5\% of gold & 2.2\% of gold \\
Step 1.5, all samplers & \textbf{none} & --- & --- \\
\bottomrule
\end{tabular}
\end{center}

On Fodors--Zagat the asymmetry is severe: the ground truth contains 110 matches, and
Step 1 supervised is trained on 50 of them. It is shown nearly half the answer key
before it selects a single attribute.

Two consequences follow. First, Step 1 supervised is better understood as an
approximate \emph{upper bound} than as a deployable baseline --- in a real
deployment, the labels it consumes are exactly what the pipeline exists to produce.
Second, the labeled sample is not held out: those 50 gold pairs remain in the
evaluation set, so a portion of the Step-1 score is measured on pairs whose
attribute selection was fitted with knowledge of their labels. On Fodors--Zagat that
overlap is 45.5\% of the ground truth; on DBLP--ACM, where the gold set is much
larger, it is 2.2\% and correspondingly less serious.

This does not make the Step-1 numbers wrong, and it does not rescue Step 1.5, which
still has to demonstrate a win on its own terms. But it reframes the comparison. The
question Step 1.5 should be judged on is not whether it beats a method trained on the
answer key --- it is how close a fully label-free method gets to one. Judged that way,
0.9222 against 0.9537 on Fodors--Zagat with zero supervision is a materially different
result from the same gap between two comparable methods.

The honest like-for-like baselines for Step 1.5 are Step 1's \emph{manual} and
\emph{heuristic} strategies, which are also label-free. Neither has current-pipeline
results on any dataset, which is a real gap in the evaluation
(Section~\ref{sec:notsettled}).

\subsection{The samplers differ in stability, not in accuracy}

Running every sampler six times on Fodors--Zagat shows that most of the apparent
ranking is noise. Among the top four, the differences in mean F1 are not
statistically distinguishable. What \emph{is} distinguishable is the spread: samplers
that draw their profile pairs randomly vary an order of magnitude more than those
that derive them deterministically.

\begin{center}
\small
\begin{tabular}{lrrrr}
\toprule
Sampler & Mean F1 & SD & TP range & Sampling \\
\midrule
diversity & 0.9222 & \textbf{0.0021} & 95--96 & deterministic \\
stratified\_diversity & 0.9152 & 0.0042 & 95--97 & deterministic \\
similarity\_stratified & 0.9110 & 0.0211 & 87--99 & stochastic \\
random & 0.8843 & 0.0429 & \textbf{79--96} & stochastic \\
importance\_based & 0.8356 & 0.0063 & 78--81 & deterministic \\
\bottomrule
\end{tabular}
\end{center}

Two things follow. First, a single run of a stochastic sampler is not evidence: one
\texttt{random} run can land anywhere from 79 to 96 true positives, spanning nearly
the full gap between the best and worst methods. Second, the pipeline's current
default, \texttt{similarity\_stratified}, is one of the unstable ones.

The one robust ranking result is negative: \texttt{importance\_based} is reliably the
worst sampler on this dataset, significantly so against all four alternatives. Since
its sampling is deterministic, this is not one unlucky draw.

\subsection{Fodors--Zagat contained a label leak}
\label{sec:leak}

The \texttt{class} attribute is the ground-truth cluster identifier --- it agrees on
110 of 110 gold pairs and is unique within each table. Step-1 supervised selection had
been choosing it, which invalidated that result entirely. All Fodors--Zagat numbers
here come from re-runs with the attribute excluded. DBLP--ACM was checked and is
clean; Amazon--Walmart's \texttt{original\_id} was checked and does not leak. The leak
was found by manual inspection; the pipeline has no automated guard.

\subsection{The selection policy breaks on wide schemas}

Step 1.5 keeps attributes until their cumulative importance reaches 0.8, capped at
five. On Amazon--Walmart's thirteen attributes this target is unreachable: using the
LLM's own importance scores, the top five carry 0.612 of the mass on average and never
more than 0.789, so reaching 0.8 would need 7.7 attributes. The cap therefore binds
for 100\% of pairs by construction, selection size is a constant 5.0, and the run
tests nothing about adaptive selection. This is a property of normalizing importance
over a wide schema, not a defect in the selector.

\section{Two Defects Found in the Measurement Apparatus}
\label{sec:defects}

Both have been fixed and covered by regression tests. Neither fix was in place for any
result above.

\subsection{Prompt field order was an uncontrolled variable}

The selector emitted its chosen subset in predicted-importance order, and that order
reached the prompt verbatim. Because predicted importance shifts with the profile
sample, reseeding permuted the prompt even when it selected exactly the same fields.
Conditioning every cross-seed comparison on DBLP--ACM:

\begin{center}
\small
\begin{tabular}{lr}
\toprule
Condition & Matcher changes its answer \\
\midrule
Same attribute set, same order & \textbf{0.10\%} \\
Same attribute set, different order & \textbf{5.44\%} \\
Different attribute set & 1.23\% \\
\bottomrule
\end{tabular}
\end{center}

Reordering the same fields changes the answer 54 times more often than re-sending an
identical prompt --- and more often than changing which fields are shown. Every
Step-1.5 run has therefore been varying two things at once, and on this dataset the
unintended one is the stronger. Both selection stages now sort into table-column
order.

\subsection{An unset token cap aborted runs while credit remained}

The LLM client never set \texttt{max\_tokens}, so every request defaulted to the
model's 65{,}535-token output ceiling. OpenRouter reserves credit against that
declared ceiling, not against actual use, and returns HTTP 402 when the balance cannot
cover it. The matcher in fact emits \emph{exactly one} completion token per call
(13{,}080 of 13{,}080 on DBLP--ACM; 2{,}665 of 2{,}665 on Fodors--Zagat). Requests were
reserving roughly four orders of magnitude more budget than they used, and nine runs
were destroyed this way. Explicit caps are now sent: 8 tokens for matching, 512 for
profiling.

A related hazard remains partly open: a failed LLM call is scored as a non-match and
the run still writes complete metrics, so a wholly failed run reports F1 $=0.0000$ and
a partly failed one reports a plausible-looking number. The sweep driver now discards
these and halts; a directly launched run does not, and must be checked.

\section{What Is Not Settled}
\label{sec:notsettled}

\textbf{The sampler ranking, on either dataset.} The DBLP--ACM figures are single runs
at one seed. Reseeding the one sampler that was successfully re-run moves F1 across
0.8847--0.9258, a spread of 0.0411 --- wider than the 0.0370--0.0539 gap the
single-seed table attributes to sampler choice. The headline result that
\texttt{importance\_based} wins on DBLP--ACM is therefore withdrawn pending repeated
seeds: not contradicted, but unsupported. Replicate noise at fixed configuration is
only 10 true positives, so the result may well survive; it has simply not been shown.

The attempted sweep completed five of thirteen runs before the budget ran out.

\textbf{Cross-generation comparability.} The canonical-order fix means no pre-fix run
may be compared against a post-fix run. The Step-1.5 tables need regenerating in full
rather than extending.

\textbf{Whether Step 1.5 beats a label-free competitor.} Every Step-1 result available
comes from the supervised strategy, which consumes ground-truth labels
(Section~\ref{sec:supervision}). The two label-free Step-1 strategies, \emph{manual}
and \emph{heuristic}, have no current-pipeline results on any dataset, so the only
comparison the evidence supports is against a supervised method --- which Step 1.5 is
not, and is not trying to be. Running those two baselines is cheap and would settle
the central question of the thesis far more directly than another sampler sweep.

\section{Status and Cost to Complete}

Work is halted for a non-technical reason: the OpenRouter account is exhausted, at
\$95.00 provisioned against \$95.25 consumed --- \$0.25 overdrawn --- and all calls
return HTTP 402. Costs below use \$0.31 per million tokens, derived from this
project's own billing (48.1\,M tokens over 46 runs in the past week against \$14.90
charged) rather than from list pricing.

\begin{center}
\small
\begin{tabular}{lrrr}
\toprule
Work item & Tokens & Est.\ cost & Runtime \\
\midrule
Step-1 \emph{manual} + \emph{heuristic}, both datasets & 9\,M & \$2.79 & 1.1\,h \\
Fodors--Zagat Step-1 baseline (re-run) & 0.7\,M & \$0.22 & 4\,min \\
Fodors--Zagat sampler sweep, $5\times5$ & 20\,M & \$6.19 & 1.7\,h \\
DBLP--ACM \texttt{full} and Step-1 baselines & 7.5\,M & \$2.32 & 55\,min \\
DBLP--ACM sampler sweep, $5\times5$ & 98\,M & \$30.35 & 8.2\,h \\
\midrule
\textbf{Total to complete the comparison} & \textbf{135\,M} & \textbf{\$41.88} & \textbf{12\,h} \\
\midrule
Amazon--Walmart, single run (blocked on policy) & 60\,M & \$18.58 & 4.3\,h \\
\bottomrule
\end{tabular}
\end{center}

Amazon--Walmart is blocked on the selection policy rather than on budget: until the
cumulative-importance threshold is made scale-free, a run on its schema saturates and
tests nothing. Spending on it first would buy another uninterpretable result.

\section{Assessment}

The pipeline works and attribute selection demonstrably helps. The uncomfortable
findings are that the simpler mechanism currently posts the better number --- a global
two-attribute subset beats per-pair adaptive selection on Fodors--Zagat, on both
accuracy and cost --- and that the sampler comparison, which was the strongest
apparent result, has shrunk under each round of scrutiny: first when a label leak was
removed, then when repeated seeds showed the ranking was inside the noise, and finally
when prompt field order proved to be a stronger driver than attribute choice.

That first finding needs qualifying rather than conceding. The two mechanisms are not
solving the same problem: Step 1's supervised strategy is trained on 50 gold matches,
which on Fodors--Zagat is 45.5\% of the entire answer key, while Step 1.5 uses no
labels at all (Section~\ref{sec:supervision}). A method given half the ground truth
should be expected to win. The result that matters is the size of the remaining gap
--- 0.9222 against 0.9537, with zero supervision --- and that gap is small enough to
be worth defending. Establishing this properly requires the two label-free Step-1
strategies, manual and heuristic, which have never been run on the current pipeline.

The rest of what survives is worth keeping. Attribute selection beats no selection by
a wide margin. One sampler is reliably bad. And the samplers separate cleanly by
stability rather than by mean accuracy, which reframes what the contribution should be
measured on --- a more defensible claim than a ranking that reseeding erases.

The immediate need is API credit, approximately \$42, to regenerate the Step-1.5
comparisons under the corrected pipeline, plus the two label-free Step-1 baselines
that would make the central comparison fair.

\end{document}
